\documentclass[11pt]{article}
\usepackage[a4paper,margin=1in]{geometry}
\usepackage[T1]{fontenc}
\usepackage{lmodern,microtype}
\usepackage{amsmath,amssymb,amsthm,mathtools}
\usepackage{booktabs,enumitem}
\usepackage{xcolor}
\usepackage[numbers,sort&compress]{natbib}
\usepackage[colorlinks=true,linkcolor=blue!55!black,citecolor=blue!55!black]{hyperref}

\newtheorem{theorem}{Theorem}[section]
\newtheorem{proposition}[theorem]{Proposition}
\newtheorem{corollary}[theorem]{Corollary}
\theoremstyle{remark}
\newtheorem{remark}[theorem]{Remark}

\title{Unit-Cap Signed-Integer Selector Obstruction}
\author{Bin Wang}
\date{July 2026}

\begin{document}
\maketitle

\begin{abstract}
RH-257 proves that integer exponents are necessary for a single-valued finite
determinant quotient.  We audit the first bounded signed lattice by assigning
each complete RH-254 shell a weight in $\{-1,0,1\}$.  Mixed-integer linear
optimization gives zero anchor passes at all 32 endpoints.  Optimal weighted
distances range from $0.106074$ to $0.353490$, which is $7.98$--$93.48$ times
the local tolerance.  The optimizations cover
$39{,}417{,}456{,}084{,}975{,}216$ lattice points in aggregate with zero
reported MIP gap.  This is a unit-cap obstruction only; larger caps and
operator realization remain open.
\end{abstract}

\section{The necessary integer lattice}

For one endpoint let $v_j\in\mathbb R^{11}$ be the order-$2$--$12$ power
vector of the $j$th conjugate-complete shell, and let $d$ be the anchored
target difference.  RH-257 shows why a continuous signed weight is not enough:
noninteger exponents create monodromy \citep{WangRH257}.  We therefore define
the first integer relaxation
\begin{equation}
 \mathcal L_1=\left\{\sum_{j=1}^Jw_jv_j:w_j\in\{-1,0,1\}\right\}.
 \label{eq:lattice}
\end{equation}

Every vector in $\mathcal L_1$ has a single-valued rational product at the
level of formal shell factors.  This remains only a necessary condition:
negative or repeated shell multiplicities must still be realized by an actual
quotient or superdeterminant.

\begin{proposition}[MILP formulation]
\label{prop:milp}
The distance
\begin{equation}
 \min_{w\in\{-1,0,1\}^J}
 \sum_{n=2}^{12}\frac{|d_n-(Vw)_n|}{n}
 \label{eq:objective}
\end{equation}
is an exact-integrality mixed-integer linear program with $J$ integer variables
and eleven nonnegative residual variables.
\end{proposition}

\begin{proof}
Introduce $u_n\ge0$ and the two linear inequalities
$d_n-(Vw)_n\le u_n$ and $(Vw)_n-d_n\le u_n$.  Minimize
$\sum_nu_n/n$ under integer bounds $-1\le w_j\le1$.
\end{proof}

\section{Finite audit}

The shell counts range from 20 to 34, so one endpoint contains $3^J$ lattice
points.  The optimizer searches this class implicitly rather than enumerating
every vector.  All 32 solves report zero MIP gap.

\begin{center}
\small
\begin{tabular}{lr}
\toprule
quantity & value\\
\midrule
endpoints & 32\\
aggregate lattice points & 39,417,456,084,975,216\\
anchor passes & 0\\
distance range & 0.106074--0.353490\\
distance/tolerance range & 7.98--93.48\\
minimum failure margin & 0.103574\\
integrality-gap range & 0.106074--0.353490\\
nonzero weights in an optimum & 20--30\\
maximum branch-and-bound nodes & 368\\
maximum reported MIP gap & 0\\
\bottomrule
\end{tabular}
\end{center}

\begin{corollary}[Scoped unit-cap obstruction]
No margin-32 complete-shell weight vector with entries in $\{-1,0,1\}$ reaches
the deterministic anchor within the archived endpoint tolerance.
\end{corollary}

\begin{proof}
At every endpoint the global MILP objective exceeds the local tolerance.  The
reported exact-integrality gap is zero.
\end{proof}

\section{Boundary and route decision}

The theorem does not exclude weights of magnitude two or larger.  Such a fit
would still not prove that the noisy operator contains the required signed
copies.  Rather than escalating an increasingly artificial multiplicity cap,
the next paper returns to the exact orthogonal quotient and its block powers,
where cancellation is operator-derived rather than fitted
\citep{WangRH245,WangRH246}.

Gates A--E remain false/open.  No Hilbert--Polya operator, Riemann-zero
identification, zeta-divisor equality, or RH implication is claimed.

\bibliographystyle{plainnat}
\bibliography{references}
\end{document}
